% ЕМТ 2023, VIII клас — точен LaTeX препис от официалните файлове в Klasirane.com.
% Условия: https://klasirane.com/api/competitions/EMT/All/2023/8%20%D0%BA%D0%BB/probs
% SHA-256 (условия): 927dfb30643ec186da524b36c38d47e54a061785107a3f8012b0c32a3c37e92e
% Решения: https://klasirane.com/api/competitions/EMT/All/2023/8%20%D0%BA%D0%BB/sol
% SHA-256 (решения): c0a52ab1609769d9bd967fd1e785a02698f3f164ec0589199f9a493846367a0a
% Точковите оценки, правописът, формулировките и видимите печатни особености са запазени от източника.
% В решението на задача 8.4 е отпечатана таблица с 10 реда и 9 стълба, въпреки че условието е за таблица 9 x 9; преписът запазва точно отпечатания пример.

Задача 8.1. Нека $a$ е най-голямата стойност на израза $24y-9y^2$, където $y$ е рационално число, а $b$ е най-малкото цяло число, изпълняващо неравенството
$$
(t+3)^3-(6t-7)^2-(t-9)^3<3.
$$
Разложете на (неразложими) множители с цели коефициенти израза
$$
a(x-1)x^3+bx-2x-1.
$$

Решение. Имаме $24y-9y^2=16-(3y-4)^2$, чиято най-голяма стойност $a=16$ се достига за $y=\dfrac43$. Даденото неравенство е еквивалентно с
$$
t^3+9t^2+27t+27-36t^2+84t-49-t^3+27t^2-243t+729<3
$$
$$
-132t+704<0,
$$
т.е. $t>\dfrac{16}{3}$ и $b=6$. Замествайки $a=16$, $b=6$ в дадения израз, получаваме
$$
16(x-1)x^3+6x-2x-1=16x^4-16x^3+4x-1=
$$
$$
=(16x^4-1)-4x(4x^2-1)=(4x^2-1)(4x^2+1)-4x(4x^2-1)=
$$
$$
=(4x^2-4x+1)(2x-1)(2x+1)=(2x-1)^3(2x+1).
$$

Оценяване. (6 точки) 2 т. за обосновано намиране на $a$; 2 т. за обосновано намиране на $b$; 2 т. за разлагане до неразложими множители и то само при правилно намерени $a,b$.

Задача 8.2. Изпъкнал четириъгълник ще наричаме иновативен, ако диагоналите му го разделят на четири триъгълника с едни и същи мерки на ъглите. Например квадратът е иновативен четириъгълник, понеже четирите триъгълника са с мерки $90^\circ$, $45^\circ$, $45^\circ$. Да се намерят мерките на ъглите на иновативен четириъгълник, ако една от тях е $13^\circ$.

Отговор. $90^\circ$, $90^\circ$, $167^\circ$, $13^\circ$ или $13^\circ$, $167^\circ$, $13^\circ$, $167^\circ$.

Решение. Нека четириъгълникът е $ABCD$ с $\angle BAD=13^\circ$ и диагоналите $AC$ и $BD$ се пресичат в $O$. Ако допуснем, че диагоналите не са перпендикулярни, то при $\angle AOB>90^\circ$ (случаят $\angle AOD>90^\circ$ е аналогичен) имаме $\angle AOB>90^\circ>\angle AOD$ и $\angle AOB>\angle OAD$, $\angle AOB>\angle ODA$ (понеже $\angle AOB$ е външен ъгъл за триъгълника $AOD$), т.е. мярката на $\angle AOB$ не се среща в триъгълника $AOD$, противоречие. Така $AC$ и $BD$ са перпендикулярни.

Нататък, ако $\angle BAO=\angle ADO$, то $\angle BAD=\angle BAO+\angle OAD=\angle BAO+90^\circ-\angle ADO=90^\circ$, противоречие с $\angle BAD=13^\circ$. Така остава само възможността $\angle BAO=\angle DAO$, като в такъв случай $AC$ разполовява $\angle BAD$, т.е. $AC$ е симетрала на $BD$. Сега от триъгълниците $AOD$ и $DOC$ следва или $\angle ADO=\angle CDO$ (в такъв случай $BD$ е симетрала на $AC$ и $ABCD$ е ромб), или $\angle ADO=\angle DCO=90^\circ-\angle CDO$, т.е. $\angle ADC=90^\circ$, аналогично $\angle ABC=90^\circ$ и четвъртият ъгъл е $\angle BCD=180^\circ-\angle BAD$.

Оценяване. (6 точки) 1 т. за верен отговор, 2 т. за доказателство, че диагоналите са перпендикулярни (не се дават точки само за предполагане на този факт), 1 т. за обосновка, че единият от диагоналите разполовява два срещуположни ъгъла на четириъгълника, по 1 т. за всеки от двата случая за другия диагонал,

Задача 8.3. Да се намерят всички двойки $(a,b)$ от взаимно прости естествени числа, такива че $a<b$ и $b$ дели
$$
(n+2)a^{n+1002}-(n+1)a^{n+1001}-na^{n+1000}
$$
за всяко естествено число $n$.

Отговор. $(a,b)=(3,5)$.

Решение. Понеже $a$ и $b$ са взаимно прости, то такива са $b$ и $a^{n+1000}$, съответно исканото е еквивалентно на $b$ да дели $(n+2)a^2-(n+1)a-n$ за всяко $n$. От $n=1$ и $n=2$ получаваме, че непременно $b$ дели $3a^2-2a-1$ и $4a^2-3a-2$. Оттук $b$ дели
$$
4(3a^2-2a-1)-3(4a^2-3a-2)=a+2
$$
и тъй като $3a^2-2a-1=3(a-2)(a+2)-2(a+2)+15$, то непременно $b$ дели 15.

Явно $b>a\ge1$, т.е. $b\ge2$. Ако $b=15$, то $b\mid a+2$ дава $a=13$, но $4\cdot13^2-3\cdot13-2=635$ не се дели на 3. Ако $b=3$, то $a=1$, но $4\cdot1^2-3\cdot1-2=-1$ не се дели на 3.

Остава $b=5$, съответно $a=3$. Действително, $(n+2)\cdot3^2-(n+1)\cdot3^1-n=5(n+3)$ се дели на 5.

Оценяване. (7 точки) 1 т. за верен отговор и проверката му; 6 т. за отхвърляне на всяка друга възможност, от които: 1 т. за свеждане до делимост на многочлени от най-много втора степен, 1 т. за фокусиране върху изрази от (поне две) малки $n\le4$, 1 т. за свеждане до делимост на два многочлена от най-много първа степен, 2 т. за извод от вида $b\mid5p$, където $p$ е просто число, 1 т. за отхвърляне на 1, $p$ и $5p$, както и на $a\ne3$ при $b=5$.

Задача 8.4. Във всяко от полетата на квадратна таблица $9\times9$ е записано цяло число. За всеки $k$ числа, намиращи се в един и същ ред (стълб), сборът им е в същия ред (стълб). Намерете най-малкия възможен брой нули в таблицата, ако:

а) $k=5$;

б) $k=8$.

Отговор. а) 63; б) 0.

Решение. а) Пример: номерираме редовете и стълбовете от 1 до 9. Записваме 1 в полетата $(i;i)$ $(i=1,\ldots,9)$; $-1$ в поле $(1;9)$ и в полетата $(i;i-1)$ $(i=2,\ldots,9)$; 0 в останалите полета. Възможните сборове са 1, 0 и $-1$.

Оценка: Да предположим, че има поне 19 ненулеви числа. От принципа на Дирихле на някой ред ще има поне три ненулеви числа, а значи и поне две ненулеви числа с еднакъв знак, да речем положителни (ситуацията при отрицателни е аналогична). Да наредим числата в този ред по големина: $a_1\le a_2\le\cdots\le a_9$, където $a_9\ge a_8>0$. Ако $a_5\ge0$, то $a_5+a_6+a_7+a_8+a_9\ge a_8+a_9>a_9$ трябва да е на същия ред: абсурд. Ако $a_5<0$, то $a_1+a_2+a_3+a_4+a_5<a_1$ трябва да е на същия ред: абсурд.

б) Възможен пример без нули е както следва (работи, понеже $5.3+3.(-4)=3$ и $4.3+4.(-4)=-4$):
$$
\begin{array}{|r|r|r|r|r|r|r|r|r|}
\hline
3&3&3&3&3&-4&-4&-4&-4\\
\hline
-4&3&3&3&3&3&-4&-4&-4\\
\hline
-4&-4&3&3&3&3&3&-4&-4\\
\hline
-4&-4&-4&3&3&3&3&3&-4\\
\hline
-4&-4&-4&-4&3&3&3&3&3\\
\hline
3&-4&-4&-4&-4&3&3&3&3\\
\hline
3&3&-4&-4&-4&-4&3&3&3\\
\hline
3&3&3&-4&-4&-4&-4&3&3\\
\hline
3&3&3&3&-4&-4&-4&-4&3\\
\hline
3&3&3&3&3&-4&-4&-4&-4\\
\hline
\end{array}
$$

Оценяване. (7 точки) а) 2 т. за работещ пример (ако проверката, че примерът работи, е неочевидна, тя трябва да присъства) и 2 т. за обоснована оценка; б) 3 т. за работещ пример (ако проверката, че примерът работи, е неочевидна, тя трябва да присъства).
